\subsection{Axes IR Declaration and Primitive Reference}
\label{app:axes-ir-declarations-primitives}
\label{app:language-declarations}
An Axes IR declaration binds ordered groups to operand positions and
specifies both output order and scalar semantics. The generalized dot below
pairs batch and contraction positions in list order; free axes preserve
operand order. Paired groups require matching extents, distinct in-range
selectors, and disjoint roles, including when groups are empty.
Transpose and slice declaration sketches illustrating the attribute roles
in Table~\ref{tab:attribute-roles} appear in Supplement~\cite{tepl-supplement},
Section~S1.

\begin{lstlisting}[language=TEPL,
  caption={Axes IR representation of generalized dot, with explicit attributes, axis bindings, and value semantics.},
  label={lst:dot-declaration}]
op dot(lhs: tensor, rhs: tensor) -> tensor {
  attrs {
    lb, lc, rb, rc: axis[];
  }
  axes {
    let lf = ordered_complement(all_axes(lhs), attrs.lb ++ attrs.lc);
    let rf = ordered_complement(all_axes(rhs), attrs.rb ++ attrs.rc);
    bind B { lhs: attrs.lb, rhs: attrs.rb };
    bind K { lhs: attrs.lc, rhs: attrs.rc };
    bind I { lhs: lf };
    bind J { rhs: rf };
    yield B ++ I ++ J;
  }
  semantics {
    yield Reduce(K, Product(lhs, rhs));
  }
}
\end{lstlisting}

\paragraph{Value and index expressions.}
\label{app:value-index-expressions}
The value primitives below compose operation semantics. Nonbroadcasting
addition binds corresponding operand axes and preserves their shared
order. Calculated reads use coordinate sequences, with one entry per
accessed axis. For example, a static slice reads at
\code{starts + strides * output_indices}. These coordinate calculations
are supported only for the modeled contracts and are outside the direct
whole-axis completeness theorem.
\begin{table}[!htbp]
\caption{Axes IR value and index primitives.}
\label{tab:axes-ir-value-index-primitives}
\small
\begin{tabularx}{\linewidth}{@{}lY@{}}\toprule
Primitive & Interpretation\\\midrule
\code{input}, \code{Read(input, indices)} & Operand read through axis bindings or explicit coordinates.\\
\code{Const(c)} & Signed integer scalar constant.\\
\code{Add(a,b)}, \code{Product(a,b)} & Scalar addition and multiplication.\\
\code{Reduce(G,e)} & Finite summation over group \code{G}; other reductions are not implied.\\
\code{Coords(G)} & Current coordinates in group order.\\
\code{IndexAdd(a,b)}, \code{Concat(a,b)} & Elementwise index addition and sequence concatenation.\\
\code{IndexScale(s,i)}, \code{IndexOffset(o,i)} & Entrywise scaling and offset of index sequence \code{i}.\\\bottomrule
\end{tabularx}
\end{table}

